\documentclass[11pt]{amsart}
\usepackage{amsmath,amssymb,amsthm}
\usepackage{graphicx}
\usepackage[hidelinks]{hyperref}

\newtheorem{theorem}{Theorem}
\newtheorem{lemma}{Lemma}
\newcommand{\cL}{\mathcal L}
\newcommand{\N}{\mathbb N}

\title[Dual-center commutants on block sums]{The dual-center commutant problem on uncountable block sums}
\author{fa\_banach\_001, agent\_lane\_01}
\date{2026-08-11}

\begin{document}
\maketitle

\section*{Status and source question}

This packet gives a partial positive answer to the open problem in Remark~4
of Orhon \cite{Orhon2010}.  If $i:Z(E)\to\cL(E)$ is the natural embedding and
$E$ has a topologically full center, the question is:
\[
 T\in\cL(E'),\qquad Ti^*(a)=i^*(a)T\quad(a\in Z(E))
 \quad\Longrightarrow\quad T\in Z(E')\ ?
\]
The cited book \cite[\S9.11(1)]{AAK} proves this when $E$ has a
quasi-interior point.  The theorem below treats arbitrary-index $c_0$- and
$\ell_p$-sums of such lattices.  For an uncountable index set, the resulting
lattice generally has no quasi-interior point.

\begin{center}
\includegraphics[width=0.96\textwidth]{figures/open_problem_crop.png}
\end{center}

\section*{Definitions and notation}

The ideal center of a Banach lattice $F$ is
\[
 Z(F)=\{S\in\cL(F): |S|\leq C I_F\text{ for some }C\geq0\}.
\]
Equivalently, it is the algebra of bounded orthomorphisms dominated in
modulus by a scalar multiple of the identity.

Let $\Gamma$ be any index set and let $(E_\gamma)_{\gamma\in\Gamma}$ be
Banach lattices.  We use the coordinatewise lattice structure on
\[
 E=\Big(\bigoplus_{\gamma\in\Gamma}E_\gamma\Big)_p
 \quad(1\leq p<\infty),
 \qquad\text{or}\qquad
 E=\Big(\bigoplus_{\gamma\in\Gamma}E_\gamma\Big)_{c_0}.
\]
In the $c_0$ case, $\|x_\gamma\|\to0$ means that for every $\varepsilon>0$
only finitely many coordinates have norm at least $\varepsilon$.

\section*{Idea of the proof}

The coordinate projections belong to $Z(E)$.  Consequently any operator on
$E'$ commuting with the adjoint copy of $Z(E)$ must be block diagonal.  Its
$\gamma$th block commutes with all adjoints of $Z(E_\gamma)$, because a
component-center operator can be extended by zero to a member of $Z(E)$.
The known quasi-interior theorem therefore makes every block central.  The
block norms are uniformly bounded by the norm of the original operator, so
the blocks reassemble into a member of $Z(E')$.

The same finite-coordinate approximation also proves directly that the
primal block sum has a topologically full center.

\section*{Partial theorem}

\begin{theorem}\label{thm:main}
Let $\Gamma$ be an arbitrary index set.  Suppose that every Banach lattice
$E_\gamma$ has a quasi-interior point, and let $E$ be either
\[
 \Big(\bigoplus_{\gamma\in\Gamma}E_\gamma\Big)_p
 \quad\text{for some }1\leq p<\infty,
 \qquad\text{or}\qquad
 \Big(\bigoplus_{\gamma\in\Gamma}E_\gamma\Big)_{c_0}.
\]
Then $E$ has a topologically full center.  Moreover, if
$T\in\cL(E')$ commutes with $i^*(Z(E))$, then $T\in Z(E')$.
\end{theorem}

\begin{proof}
For each $\gamma$, let $Q_\gamma$ be the coordinate band projection on $E$.
The standard coordinate argument gives the isometric algebra and lattice
identification
\begin{equation}\label{eq:center}
 Z(E)=\Big\{(S_\gamma)_{\gamma\in\Gamma}:
 S_\gamma\in Z(E_\gamma),\ 
 \sup_\gamma\|S_\gamma\|<\infty\Big\},
\end{equation}
where the family acts coordinatewise.  Indeed, a central operator preserves
the bands $Q_\gamma E$ and its restrictions are central; the converse follows
immediately from $|S_\gamma|\leq\|S_\gamma\|I_{E_\gamma}$.

We first verify topological fullness.  Fix $x=(x_\gamma)\in E_+$ and put
\[
 J_\gamma=\overline{Z(E_\gamma)x_\gamma}\subset E_\gamma.
\]
Because $E_\gamma$ has a quasi-interior point, its center is topologically
full \cite{Orhon2010}, so $J_\gamma$ is an ideal.  We claim that
\begin{equation}\label{eq:orbit}
 \overline{Z(E)x}=
 \Big(\bigoplus_{\gamma\in\Gamma}J_\gamma\Big)_p
 \quad\text{or}\quad
 \Big(\bigoplus_{\gamma\in\Gamma}J_\gamma\Big)_{c_0},
\end{equation}
respectively.  One inclusion follows from \eqref{eq:center}.  For the other,
take $y$ in the right-hand side and $\varepsilon>0$.  Truncate $y$ to a
finite set $F\subset\Gamma$ with error below $\varepsilon/2$.  For every
$\gamma\in F$, choose $S_\gamma\in Z(E_\gamma)$ so that $S_\gamma x_\gamma$
approximates $y_\gamma$ closely enough that the finite-sum error is below
$\varepsilon/2$.  Extending this finite family by zero gives an element of
$Z(E)$ and proves \eqref{eq:orbit}.  Its right-hand side is an ideal, proving
that $E$ has a topologically full center.

Now suppose that $T\in\cL(E')$ commutes with $i^*(Z(E))$.  Since
$Q_\gamma\in Z(E)$, we have
\[
 TQ_\gamma'=Q_\gamma'T\qquad(\gamma\in\Gamma).
\]
Thus $T$ leaves $Q_\gamma'E'\cong E_\gamma'$ invariant; denote the
restriction by $T_\gamma$.  If $S\in Z(E_\gamma)$, extend $S$ by zero on all
other coordinates.  Formula \eqref{eq:center} puts the extension in $Z(E)$,
so $T_\gamma$ commutes with the natural adjoint copy of $Z(E_\gamma)$.
The quasi-interior theorem \cite[\S9.11(1)]{AAK} yields
\begin{equation}\label{eq:blockcenter}
 T_\gamma\in Z(E_\gamma')
 \qquad\text{and}\qquad
 \|T_\gamma\|\leq\|T\|.
\end{equation}

Let $q$ be conjugate to $p$.  The dual is
\[
 E'=\Big(\bigoplus_{\gamma\in\Gamma}E_\gamma'\Big)_q
\]
with the usual interpretation $q=\infty$ when $p=1$; in the $c_0$ case it is
the $\ell_1$-sum.  For every $x'\in E'$ and every $\gamma$,
\[
 (Tx')_\gamma=Q_\gamma'Tx'=TQ_\gamma'x'=T_\gamma x_\gamma'.
\]
This coordinate identity also handles $q=\infty$: an element of the product
is determined by all its coordinates, so there is no additional residual
operator.  By \eqref{eq:blockcenter}, the uniformly bounded family
$(T_\gamma)$ defines a coordinatewise central operator on $E'$, dominated in
modulus by $\|T\|I_{E'}$.  It agrees with $T$ on every coordinate and hence
equals $T$.  Therefore $T\in Z(E')$.
\end{proof}

\section*{Why the full problem remains}

Orhon's topological-fullness theorem identifies
$Z(E')=i^*(Z(E)'')$ through an Arens extension.  Approximating an element of
$Z(E')$ by elements of $i^*(Z(E))$ gives only weak-star operator convergence.
An arbitrary $T\in\cL(E')$ need not be weak-star continuous, so its
commutation cannot be passed through that limit.  The block-sum proof avoids
this by using actual coordinate projections already in $Z(E)$.

The method does not automatically cover $\ell_\infty$-sums.  Their Banach
duals contain a singular part beyond the coordinate dual bands, and the
coordinate equation no longer exhausts the dual.  Controlling or exploiting
that singular part is a plausible route to the unrestricted problem.

\section*{Verification and novelty notes}

The proof is structural and uses no numerical computation.  The primal-center
identification, topological-fullness argument, dual coordinate identity, and
$p=1$ product case were checked separately.  The local run indexes contained
no duplicate.  A local source search found Alpay--Orhon \cite{AlpayOrhon2014},
which proves a related theorem for adjoints of order-bounded primal operators;
that extra adjoint continuity does not apply to arbitrary operators on $E'$.
Bounded exact-phrase web searches on 2026-08-11 found no direct later
resolution, but some final targeted calls timed out, so novelty confidence is
moderate.  Human review should focus on the standard center identifications
and the cited quasi-interior theorem.

\begin{thebibliography}{3}

\bibitem{AAK}
Y.~A. Abramovich, E.~L. Arenson, and A.~K. Kitover,
\emph{Banach $C(K)$-modules and operators preserving disjointness},
Pitman Research Notes in Mathematics Series 277, Wiley, 1992.

\bibitem{AlpayOrhon2014}
S. Alpay and M. Orhon,
\emph{Characterization of Riesz spaces with topologically full center},
Positivity 19 (2015), 799--816; arXiv:1406.6335.

\bibitem{Orhon2010}
M. Orhon,
\emph{The ideal center of the dual of a Banach lattice},
Positivity 15 (2011), 555--564; arXiv:1002.4346.

\end{thebibliography}

\end{document}
